\scanpage{15}
on obtient des bijections inverses l'une de l'autre
\diagram{\begin{tikzcd}[ampersand replacement=\&]
\pi\arrow[r,shift left,"\exp"] \& \gamma\arrow[l,shift left,"\log"]
\end{tikzcd}}

Bien entendu, $\pi$ est un idéal \inserted{(nilpotent)} de $L_\Lambda(M)$, et
à fortiori une sous-algèbre de Lie. En
transportant à $\pi$ la structure de groupe de
$\gamma$, on trouve une structure de groupe dans
$\pi$, qui se déduit de sa structure
d'algèbre de Lie par la formule de
Hausdorff-Campbell. Si $M$ est une
$\Lambda$-algèbre \inserted{$\deleted{A}B$}, on trouve le
\uline{corollaire} les \uline{automorphismes} de la
$\Lambda$-algèbre $\underline{\deleted{A}B}$ (où $\Lambda$ est une algèbre
$k$-augmentée à \deleted{\unsure{noyau}} idéal d'augmentation
nilpotent, $k$ étant une $\mathbb Q$-algèbre)
qui induisent l'identité sur $\deleted{\unsure{A}}B\otimes_\Lambda k$
correspondent biunivoquement aux
$\Lambda$-dérivations \inserted{$D$} de $B$ telles que $DB\subset I(\Lambda)B$.
% archive sheet number: 14

